\begin{abstract}
Hybrid classical-quantum machine learning (QML) models for medical image analysis frequently confront three fundamental bottlenecks: classical gradient collapse where gating routers degenerate into purely classical pathways, the spatial reconstruction paradox where severe quantum latent constraints destroy high-resolution anatomical boundaries, and the absence of parameter-matched classical controls. 
In this work, we present Q-FUSION, an adaptive classical-quantum architecture designed to operate under strict latent-dimension and variational circuit constraints. 
To resolve classical collapse, we introduce an annealed regularized routing objective that balances early quantum-classical pathway exploration before converging to sample-dependent allocation. 
To resolve the spatial reconstruction paradox, we formulate quantum-conditioned spatial modulation via Feature-wise Linear Modulation (FiLM), allowing compact 8-dimensional quantum state measurements to modulate high-resolution 3D convolutional feature maps without high-dimensional unflattening layers. 
To eliminate parameter count confounders, all quantum variational circuits are paired with capacity-matched classical surrogate transformations achieving strict 0.0000\% capacity variance.
We evaluate Q-FUSION across two cohorts: 150 patients from the multi-sequence BraTS2021 benchmark across one to four physical imaging sequences (T1, T1ce, T2, FLAIR) and 181 patients from the SynthRAD2023 cross-modality benchmark (MRI and CT). 
On BraTS2021, when input complexity expands to full quad-modal imaging ($M=4$), classical baseline pooling degrades by 4.41\% in Dice Similarity Coefficient (DSC) due to cross-modal feature dilution, whereas adaptive routing maintains 0.7015 DSC (+4.79\% over the classical baseline) through dynamic representation reallocation ($\bar{\alpha} = 0.4883$). 
Under multi-seed evaluation on SynthRAD2023, adaptive routing achieves a +2.47\% mean DSC advantage over rigid quantum fusion ($0.8039 \pm 0.0945$ vs. $0.7792 \pm 0.1011$). 
Furthermore, under simulated physical gate depolarizing noise ($p \in [0.000, 0.050]$), adaptive routing limits zero-shot degradation to 0.01\% DSC compared to a monotonic 0.32\% loss in rigid quantum architectures. 
Our findings demonstrate that adaptive classical-quantum routing provides architectural resilience against cross-modal interference and NISQ hardware noise under identical resource constraints.
\end{abstract}
